\documentclass[11pt,a4paper]{article}
\usepackage{curso}
\begin{document}
\portada{Lección}{4.10}{Optimización con Ising y lógica probabilística invertible}

\noindent\textbf{Unidad 4.}\quad\textbf{Objetivos:} Max-Cut y QUBO como estados fundamentales de Ising; compuertas como penalizaciones; lógica invertible y
factorización fijando un multiplicador; límites del enfoque.
\textbf{Prerrequisitos:} 4.4 (recocido), 4.5 (pbits).

\section{Optimización como estado fundamental}
A $T\to0$, $p\propto e^{-\beta E}$ se concentra en los mínimos de $E$. Si la solución de un problema es el mínimo de una energía de Ising, un chip de
pbits la busca enfriándose \cite{lucas}.
\begin{resultado}[ (Max-Cut)]
Con $s_i=\pm1$ indicando el lado de cada nodo, $\text{corte}(s)=\sum_{i<j}w_{ij}(1-s_is_j)/2$. Maximizar el corte equivale a minimizar $E=-\sum J_{ij}s_is_j$
con $J_{ij}=-w_{ij}$ (antiferromagnético).
\end{resultado}
\begin{ejemplo}
Grafo aleatorio con $n=10$ nodos y 18 aristas: corte máximo 14 (fuerza bruta, $2^{10}$ estados). Recocido con pbits ($T:2\to0{,}05$, 300 barridos):
corte 14 en 100 de 100 corridas \evS.
\end{ejemplo}
QUBO: variables $x\in\{0,1\}$ con costo cuadrático; $x=(1+s)/2$ lo convierte en Ising (más una constante). SAT: una penalización por cláusula; la energía mínima
es 0 si y solo si la fórmula es satisfacible.

\section{Compuertas como energías y lógica invertible}
\begin{definicion}[ (penalización AND)]
$P_{\text{AND}}(u,v,z)=uv-2uz-2vz+3z$, con $u,v,z\in\{0,1\}$, vale 0 exactamente en las cuatro filas de $z=u\wedge v$ y $\ge1$ en las demás.
\end{definicion}
Un circuito es la suma de las penalizaciones de sus compuertas (compartiendo variables). A $\beta=2$ las cuatro filas válidas tienen probabilidad $0{,}227$
cada una. Fijando la salida:
\begin{center}\begin{tabular}{lcccc}\toprule
$(u,v)$ & 00 & 01 & 10 & 11\\\midrule
$p(u,v\mid z=1)$ & 0,002 & 0,106 & 0,106 & 0,785\\
$p(u,v\mid z=0)$ & 0,319 & 0,319 & 0,319 & 0,043\\\bottomrule\end{tabular}\end{center}
El mismo circuito calcula hacia adelante (entradas fijas) y hacia atrás (salidas fijas): lógica probabilística invertible \cite{camsari}.

\section{Factorizar fijando un multiplicador}
Fijar la salida de un multiplicador en $N$ hace que las entradas busquen $p,q$ con $pq=N$. Modelo de juguete: $N=143$, $p,q$ impares de 4 bits, energía
$(N-pq)^2$, recocido de Metropolis: $11\times13$ en 79 de 200 corridas. No rompe la criptografía: el paisaje es rugoso, $(N-pq)^2$ tiene términos de orden
alto (requiere espines auxiliares) y el tiempo crece rápidamente con el número de bits. La Unidad 7 presenta un algoritmo por difusión con garantías.

\begin{aplicacion}
Programación de tareas, rutas, carteras (QUBO); p-computadores en FPGA y máquinas de Ising dispersas \cite{aadit}.
\end{aplicacion}

\begin{error}
\begin{itemize}[nosep]
\item Penalizaciones débiles frente al objetivo. \quad Olvidar la constante QUBO $\leftrightarrow$ Ising.
\item Esperar tiempo polinomial en problemas NP-difíciles. \quad Términos de orden alto sin espines auxiliares.
\end{itemize}
\end{error}

\begin{thebibliography}{9}\small
\bibitem[17]{camsari} K. Y. Camsari, B. M. Sutton y S. Datta, p-bits for probabilistic spin logic, Appl. Phys. Rev. 6, 011305 (2019).
\bibitem[19]{aadit} N. A. Aadit et al., Massively parallel probabilistic computing with sparse Ising machines, Nat. Electron. 5, 460 (2022).
\bibitem{lucas} A. Lucas, Ising formulations of many NP problems, Front. Phys. 2, 5 (2014).
\end{thebibliography}
\end{document}
